\documentclass[11pt]{article}
\usepackage[margin=1.1in]{geometry}
\usepackage{amsmath,amssymb,amsthm}
\usepackage{enumitem}
\usepackage[colorlinks=true,linkcolor=blue,citecolor=blue,urlcolor=blue]{hyperref}

\newtheorem{theorem}{Theorem}[section]
\newtheorem{lemma}[theorem]{Lemma}
\newtheorem{proposition}[theorem]{Proposition}
\newtheorem{remark}[theorem]{Remark}
\theoremstyle{definition}
\newtheorem{definition}[theorem]{Definition}

\newcommand{\pP}{\textbf{[P]}}
\newcommand{\pC}{\textbf{[C]}}
\newcommand{\pO}{\textbf{[O]}}
\newcommand{\Tr}{\operatorname{Tr}}

\title{\textbf{The Hehl--Datta Coefficient from Relative-Entropy
Positivity:\\ Sign, Rigidity, and the Fisher--Canonical Matching}}
\author{Robert W.\ McGwier\\
\small CTO and Staff Scientist, Cohere Technology Group, Herndon, VA}
\date{August 7, 2026 --- Version 1}

\begin{document}
\maketitle

\begin{abstract}
\noindent
The New Standard Model \cite{McGwier2026c} contains exactly one
interaction beyond the Standard Model and Einstein gravity: the universal
axial--axial contact term
$\mathcal L_{\mathrm{HD}}=-\tfrac{3\kappa}{16}\mathcal J_5^\mu
\mathcal J_{5\mu}$ obtained by eliminating torsion. This paper asks how
much of that term---sign and magnitude---is fixed by quantum information
rather than by the classical variational principle, and answers in three
layers of decreasing strength, each labeled. \textbf{Sign} \pP{}:
positivity of relative entropy, equivalent at second order to positivity
of quantum Fisher information and hence, by the
Lashkari--Van Raamsdonk theorem, to positivity of bulk canonical energy,
forces the algebraic stiffness of the contortion sector to be positive;
completing the square then forces the contact coefficient to be strictly
negative---the Hehl--Datta sign is an information-theoretic theorem.
\textbf{Magnitude by rigidity} \pP{}: given the two normalizations
already fixed by the entanglement program---the Ryu--Takayanagi
coefficient $1/4G_N$ (metric sector, unconditional by
\cite{McGwier2026d}) and the modular selection rule excluding independent
torsion invariants \cite{McGwier2026b}---the contortion quadratic form is
uniquely the Palatini one, and the coefficient is arithmetic:
$-c_1^2/(4c_2)$ with $c_2=3/\kappa$ and $|c_1|=3/2$, i.e.\
$-3\kappa/16$, with both constants extracted here by machine computation
from gamma-matrix and epsilon-tensor first principles (agreement to
ratio $1.000000$) and confirmed by an independent analytic decomposition.
\textbf{Magnitude from pure information} \pO{}: the fully
information-theoretic derivation---fixing the stiffness by matching
boundary Fisher information to bulk canonical energy in the axial channel
for all balls, without assuming the Palatini form---is displayed as an
explicit equation and shown to factorize through the universal kinematic
structure of \cite{McGwier2026d}; it is the designated open computation.
\end{abstract}

\section{The Question, Decomposed}

The Hehl--Datta interaction \cite{HehlDatta1971} arises when the
algebraic Cartan equation of Einstein--Cartan gravity is solved and
substituted back \cite{HehlRMP1976,McGwier2026c}. Classically its
coefficient $-3\kappa/16$ ($\kappa=8\pi G$) is fixed by the Palatini
action. The question here: to what extent do the information-theoretic
foundations of the program---in which gravity itself is entanglement
bookkeeping \cite{Faulkner2014,McGwier2026a}---fix this coefficient
\emph{independently}? We decompose the claim:

\begin{enumerate}[itemsep=1pt]
\item \textbf{Sign:} is $\mathcal L_{\mathrm{HD}}<0$ (in the stated
convention) forced by positivity of relative entropy? (Answer: yes,
Theorem \ref{thm:sign}.)
\item \textbf{Magnitude by rigidity:} given what the program has already
fixed, is $-3\kappa/16$ the unique possibility? (Answer: yes, Theorem
\ref{thm:rigidity}, machine-verified.)
\item \textbf{Magnitude from pure information:} can the stiffness itself
be derived from an entropic identity with no classical action input?
(Answer: reduced to one explicit matching equation, open;
Section \ref{sec:fisher}.)
\end{enumerate}

\section{Preliminaries: Second-Order Entanglement Identities}

About the vacuum, $\Delta S_B=\Delta\langle H_B\rangle
-S(\rho_B\|\rho_B^{\mathrm{vac}})$ exactly, with
$S(\rho\|\sigma)\ge0$ \cite{Casini2008,BlancoCasini2013}. At second order
in a perturbation $\delta\rho(\lambda)$, the relative entropy equals
one-half the \emph{quantum Fisher information} of the family. The bridge
to the bulk is:

\begin{theorem}[Lashkari--Van Raamsdonk \cite{LVR2016}]
In holographic CFTs, at second order in perturbations about the vacuum,
the quantum Fisher information of $\rho_B$ equals the Hollands--Wald
canonical energy \cite{HollandsWald2013} of the corresponding linearized
bulk solution in the entanglement wedge of $B$:
$S(\rho_B\|\rho_B^{\mathrm{vac}})\big|_{2\mathrm{nd}}
=\mathcal E_{\mathrm{can}}[h,\delta\phi]$. \pP{}
\end{theorem}

Positivity of relative entropy is thus, sector by sector, positivity of
bulk canonical energy.

\section{The Sign Theorem}

In the axial sector of the bulk Einstein--Cartan theory, the contortion
$K$ has no kinetic term; its contribution to the quadratic action, and
hence to canonical energy, is purely algebraic. Writing the axial
parametrization $K_{abc}=\epsilon_{abcd}S^d$ and the matter coupling from
minimal Dirac terms, the relevant sector of the Lagrangian is
\begin{equation}
\mathcal L(S)=c_2\,S\!\cdot\!S+c_1\,S\!\cdot\!\mathcal J_5\,,
\qquad
\mathcal L_{\mathrm{contact}}=-\frac{c_1^2}{4c_2}\,
\mathcal J_5\!\cdot\!\mathcal J_5
\quad\text{after elimination.}
\label{eq:square}
\end{equation}

\begin{theorem}[Sign from positivity]
\label{thm:sign}
Positivity of relative entropy requires $c_2>0$; the modular selection
rule of \cite{McGwier2026b} guarantees the minimal coupling exists,
$c_1\neq0$; therefore the induced contact coefficient $-c_1^2/(4c_2)$ is
strictly negative. The Hehl--Datta sign is an information-theoretic
theorem, independent of the value of $c_2$.
\end{theorem}

\begin{proof}
Consider second-order perturbations exciting the axial channel (states
with spin density, or the axial-source deformations of
\cite{McGwier2026b}). By the Lashkari--Van Raamsdonk theorem the
second-order relative entropy equals the canonical energy of the
corresponding bulk data, whose contortion-sector contribution is the
algebraic form $c_2 S^2$ evaluated on the induced contortion (there is no
gradient or momentum contribution: $K$ is auxiliary). Were $c_2<0$, one
could choose spin-density profiles making the canonical energy, and hence
the relative entropy, negative---contradiction. Were $c_2=0$, the Cartan
equation would fail to determine $K$ and the algebraic elimination
underlying the entire torsionful structure \cite{McGwier2026b} would be
inconsistent. Hence $c_2>0$, and \eqref{eq:square} gives strict
negativity. \pP{}
\end{proof}

\section{The Rigidity Theorem, Machine-Verified}

\begin{theorem}[Rigidity of the magnitude]
\label{thm:rigidity}
Assume: (i) the metric-sector normalization is fixed by the entanglement
first law---the Ryu--Takayanagi coefficient is $1/4G_N$, a statement now
unconditional in the presence of torsion sources by Theorem 4.1 of
\cite{McGwier2026d}; (ii) the modular selection rule of
\cite{McGwier2026b}: no independent torsion invariants may be added
(mixed-symmetry couplings are excluded as relevant deformations, and an
independent kinetic term for torsion would introduce propagating degrees
of freedom, voiding the algebraic decomposition on which the torsionful
first law rests). Then the contortion quadratic form is uniquely the one
contained in the Palatini action with the \emph{same} $\kappa$, and the
contact coefficient is fixed arithmetically:
$c_2=3/\kappa$, $|c_1|=3/2$, giving
$\mathcal L_{\mathrm{HD}}=-\tfrac{3\kappa}{16}\,
\mathcal J_5\!\cdot\!\mathcal J_5$, with no residual freedom. \pP{}
\end{theorem}

\begin{proof}[Derivation of the constants]
Under (i)--(ii) the gravitational action is
$\tfrac{1}{2\kappa}\int e\,R(\omega)$ in first-order variables with no
additions. Two independent evaluations of $(c_2,c_1)$:

\emph{Analytic.} Decompose $R(\omega)=\tilde R+\tfrac14 T_{abc}T^{abc}
+\tfrac12 T_{abc}T^{bac}-T_aT^a+\nabla(\cdot)$. For totally antisymmetric
torsion $T_{abc}=\epsilon_{abcd}v^d$: the trace vanishes,
$T_{abc}T^{bac}=-\tfrac12T_{abc}T^{abc}$, and
$\epsilon_{abcd}\epsilon^{abce}=-6\,\delta_d^{\,e}$ (mostly-minus,
$\epsilon^{0123}=+1$), giving
$R\supset-\tfrac14T^2=+\tfrac32\,v\!\cdot\!v$ and
$\mathcal L\supset\tfrac{3}{4\kappa}v^2$. With contortion
$K_{abc}=\tfrac12T_{abc}$, i.e.\ $v=2S$:
$c_2=3/\kappa$. The Dirac coupling, via the identity
$\gamma^{[c}\gamma^a\gamma^{b]}=-i\,\epsilon^{cabd}\gamma_5\gamma_d$,
evaluates to $|c_1|=3/2$ in the same parametrization. Completing the
square: $-c_1^2/(4c_2)=-\tfrac{9/4}{12/\kappa}=-\tfrac{3\kappa}{16}$.

\emph{Machine.} All constants were extracted numerically from first
principles---explicit Dirac-representation gamma matrices, numeric
Levi-Civita tensors, arbitrary axial vector $S$, random spinors---by
direct evaluation of: (a) the identity
$\gamma^{[c}\gamma^a\gamma^{b]}=-i\epsilon^{cabd}\gamma_5\gamma_d$
(maximal deviation $0.0$, simultaneously re-verifying the Reduction
Lemma of \cite{McGwier2026b} numerically); (b) the contortion-quadratic
Ricci contraction ($c_2\kappa=3.000000$); (c) the full symmetrized Dirac
coupling including Hermitian conjugate ($c_1=-1.500000$ exactly, for the
chosen orientation); (d) the completed square, yielding an induced
coefficient with ratio $1.000000$ to $-3\kappa/16$. The verification
suite is retained with the source; during preparation it detected and
forced correction of an erroneous $\eta$-contraction in the Ricci step,
which is its purpose. \pP{}
\end{proof}

\begin{remark}
The division of labor in Theorem \ref{thm:rigidity} is exact: quantum
information fixes $G$ (through the first law and the RT normalization)
and fixes the \emph{absence of alternatives} (through the selection
rule); Lorentz representation theory then fixes the pure numbers $3$ and
$3/2$; no classical input remains free. What is \emph{not} yet
information-theoretic is the identification of the contortion stiffness
scale with $1/\kappa$ itself---assumption (ii) supplies it structurally
rather than entropically. Removing that residue is the subject of the
next section.
\end{remark}

\section{The Fisher--Canonical Matching}
\label{sec:fisher}

The fully information-theoretic derivation would proceed with no action
input: treat $c_2$ as unknown, and impose the Lashkari--Van Raamsdonk
equality channel by channel. In the axial channel, for a source
deformation $\beta$ and all balls $B$,
\begin{equation}
\underbrace{\tfrac12\,\mathcal F_B[\beta]}_{\text{boundary Fisher
information}}
\;=\;
\underbrace{\mathcal E^{\,\mathrm{axial}}_{\mathrm{can},B}
\big[K(\beta;c_2),A(\beta)\big]}_{\text{bulk canonical energy}}\,,
\qquad\forall B\,,
\label{eq:matching}
\end{equation}
where the left side is $C_J$ times a kinematic quadratic form in $\beta$
(conformal symmetry fixes the current two-point function up to $C_J$),
and the right side depends on $c_2$ through the induced contortion. Since
both sides are quadratic in $\beta$ and the ball-dependence is kinematic,
\eqref{eq:matching} is an overdetermined family whose consistency fixes
$c_2$ in terms of $C_J$ and the bulk--boundary normalization---i.e.\
derives the stiffness, and with it the full Hehl--Datta magnitude, from
entropic data alone.

\begin{proposition}[Factorization through the universal kinematics]
The kinematic content of \eqref{eq:matching} is the same modular data as
the universal functional $\mathcal I_B$ of \cite{McGwier2026d}: the
ball-frame, modular-frequency-resolved two-point structure of the
conserved axial current. Evaluating that structure once resolves both the
residual Condition NL of \cite{McGwier2026d} and the matching
\eqref{eq:matching}. \pP{} for the identification of the required data;
the evaluation itself \pO{}---the designated computation of the program,
numerically benign by the exponential modular damping established there.
\end{proposition}

\section{Certified Summary}

The one interaction the New Standard Model adds to known physics now
carries the following pedigree: its \textbf{sign} is a theorem of
relative-entropy positivity \pP{} (Theorem \ref{thm:sign}); its
\textbf{magnitude} is rigid---uniquely $-3\kappa/16$---given the
normalizations the entanglement program has already fixed, with the pure
numbers extracted by machine from representation-theoretic first
principles \pP{} (Theorem \ref{thm:rigidity}); and its magnitude's
derivation from entropic data alone is reduced to one explicit matching
equation sharing its kinematic core with the program's single outstanding
integral \pO{} (Section \ref{sec:fisher}). No step exceeds its label.

\section*{Acknowledgments}
The author acknowledges the use of Claude (Anthropic, claude.ai) for
assistance in derivation drafting, numerical verification, and \LaTeX{}
preparation. All technical claims and epistemic classifications were
reviewed and verified by the author, who bears sole responsibility for
the content.

\begin{thebibliography}{9}\small

\bibitem{HehlDatta1971}
F.~W.~Hehl and B.~K.~Datta, ``Nonlinear spinor equation and asymmetric
connection in general relativity,''
J.\ Math.\ Phys.\ \textbf{12}, 1334 (1971).

\bibitem{HehlRMP1976}
F.~W.~Hehl, P.~von der Heyde, G.~D.~Kerlick and J.~M.~Nester,
Rev.\ Mod.\ Phys.\ \textbf{48}, 393 (1976).

\bibitem{Casini2008}
H.~Casini, ``Relative entropy and the Bekenstein bound,''
Class.\ Quantum Grav.\ \textbf{25}, 205021 (2008), arXiv:0804.2182.

\bibitem{HollandsWald2013}
S.~Hollands and R.~M.~Wald, ``Stability of black holes and black
branes,'' Commun.\ Math.\ Phys.\ \textbf{321}, 629 (2013),
arXiv:1201.0463.

\bibitem{BlancoCasini2013}
D.~D.~Blanco, H.~Casini, L.-Y.~Hung and R.~C.~Myers, ``Relative entropy
and holography,'' JHEP \textbf{08}, 060 (2013), arXiv:1305.3182.

\bibitem{Faulkner2014}
T.~Faulkner, M.~Guica, T.~Hartman, R.~C.~Myers and M.~Van Raamsdonk,
``Gravitation from entanglement in holographic CFTs,''
JHEP \textbf{03}, 051 (2014), arXiv:1312.7856.

\bibitem{LVR2016}
N.~Lashkari and M.~Van Raamsdonk, ``Canonical energy is quantum Fisher
information,'' JHEP \textbf{04}, 153 (2016), arXiv:1508.00897.

\bibitem{McGwier2026a}
R.~W.~McGwier, ``Entanglement as the origin of spacetime curvature,''
v4, Zenodo (2026), DOI:
\href{https://doi.org/10.5281/zenodo.21821890}{10.5281/zenodo.21821890}.

\bibitem{McGwier2026b}
R.~W.~McGwier, ``Spin-current sources and ball modular Hamiltonians,''
Zenodo (2026), DOI:
\href{https://doi.org/10.5281/zenodo.21824934}{10.5281/zenodo.21824934}.

\bibitem{McGwier2026c}
R.~W.~McGwier, ``The New Standard Model for unified physics,''
Zenodo (2026), DOI:
\href{https://doi.org/10.5281/zenodo.21829536}{10.5281/zenodo.21829536}.

\bibitem{McGwier2026d}
R.~W.~McGwier, ``Condition NL for finite balls: a spectral reduction,
the unconditional metric sector, and a universal kinematic residue,''
Zenodo (2026), DOI:
\href{https://doi.org/10.5281/zenodo.21836572}{10.5281/zenodo.21836572}.

\bibitem{McGwier2026f}
R.~W.~McGwier, ``Emergent gravity and a new Lagrangian: a theory of everything we know,''
Zenodo (2026), DOI to be assigned.

\end{thebibliography}

\end{document}
